\chapter{Workbook Phần 5: 20 Tình Huống Phân Tích Nghiệp Vụ \& Đo Lường Hiệu Quả (Business Analytics \& Metrics)}

\begin{lythuyetbox}[Mục Tiêu Luyện Tập Phần 5]
Phần này cung cấp \textbf{20 tình huống và bài toán kinh doanh định lượng thực tế} mà mọi Data Analyst, Product Analyst và Business Intelligence Engineer phải đối mặt hàng ngày. Các bài tập rèn luyện kỹ năng phân tích chỉ số tài chính, phễu chuyển đổi, định giá, phân khúc khách hàng, và cây chỉ số nghiệp vụ (KPI Trees).
\end{lythuyetbox}

\begin{baitapbox}[Hệ Thống 20 Bài Toán Nghiệp Vụ \& Phân Tích Chỉ Số]
\begin{enumerate}
    \item \textbf{Bài 1 (Tính Customer Lifetime Value -- CLV)}: Một sàn thương mại điện tử có Giá trị đơn hàng trung bình $AOV = \$85$, Tần suất mua hàng trung bình $f = 3.2$ lần/năm, Biên lợi nhuận gộp $Gross\ Margin = 40\%$, và Tỷ lệ rời bỏ khách hàng hàng năm $Churn\ Rate = 20\%$. Hãy tính Giá trị vòng đời khách hàng ($CLV$) theo mô hình đơn giản và phân tích tác động nếu Churn Rate giảm xuống còn $15\%$.
    \item \textbf{Bài 2 (CAC \& Thời Gian Hoàn Vốn Payback Period)}: Trong Quý 1, tổng chi phí Marketing và Sales là $\$120,000$. Công ty mang về 3,000 khách hàng mới, trong đó 1,000 khách đến từ kênh tự nhiên (Organic) và 2,000 khách đến từ quảng cáo trả phí (Paid). Doanh thu trung bình trên mỗi khách hàng hàng tháng $ARPU = \$25$, biên lợi nhuận gộp $80\%$. Tính Blended CAC, Paid CAC, và thời gian hoàn vốn (CAC Payback Period) theo tháng.
    \item \textbf{Bài 3 (Phân Tích Cohort Retention \& Churn)}: Nhóm khách hàng đăng ký tháng 01/2026 gồm 1,000 người. Số lượng người dùng còn hoạt động qua các tháng lần lượt là: Tháng 1 (1,000), Tháng 2 (650), Tháng 3 (520), Tháng 4 (480), Tháng 5 (460). Hãy tính Retention Rate từ M0 đến M4. Nhận xét về độ dốc của đường cong lưu giữ và xác định thời điểm sản phẩm đạt trạng thái ổn định (Retention Curve Flattening).
    \item \textbf{Bài 4 (Tối Ưu Hóa Phễu Chuyển Đổi -- Funnel Drop-off)}: Dữ liệu phễu thanh toán của ứng dụng bán lẻ: Lượt xem sản phẩm: 200,000 $\rightarrow$ Thêm vào giỏ hàng: 50,000 $\rightarrow$ Bấm thanh toán: 15,000 $\rightarrow$ Điền thông tin giao hàng: 12,000 $\rightarrow$ Thanh toán thành công: 3,600. Tính Conversion Rate của toàn phễu (End-to-End CR) và xác định bước có tỷ lệ thất thoát (Drop-off Rate) nghiêm trọng nhất cần ưu tiên cải tiến UX.
    \item \textbf{Bài 5 (Nghịch Lý Simpson Trong Chiến Dịch Tiếp Thị)}: Đội Marketing chạy 2 chiến dịch (Campaign A và Campaign B) trên 2 nền tảng Facebook và Google. Dữ liệu ghi nhận:
    \begin{itemize}
        \item Facebook: Campaign A có $CR = 2.0\%$ (20/1,000 clicks); Campaign B có $CR = 2.5\%$ (100/4,000 clicks).
        \item Google: Campaign A có $CR = 8.0\%$ (320/4,000 clicks); Campaign B có $CR = 8.5\%$ (85/1,000 clicks).
    \end{itemize}
    Tính tổng Conversion Rate của từng chiến dịch gộp cả hai nền tảng. Giải thích tại sao Campaign B thắng ở cả hai nền tảng riêng lẻ nhưng khi gộp lại Campaign A lại có tỷ lệ chuyển đổi cao hơn gấp bội.
    \item \textbf{Bài 6 (Cây Chỉ Số Nghiệp Vụ -- KPI Tree Decomposition)}: Doanh thu tháng này giảm $12\%$ so với tháng trước. Ban giám đốc yêu cầu phân rã nguyên nhân dựa trên công thức: $\text{Doanh Thu} = \text{Traffic} \times \text{Conversion Rate} \times \text{AOV}$. Biết rằng Traffic tăng $10\%$, AOV tăng $5\%$. Hãy tính xem Conversion Rate đã thay đổi bao nhiêu phần trăm và đưa ra kết luận điều tra nguyên nhân.
    \item \textbf{Bài 7 (Độ Co Giãn Của Cầu Theo Giá -- Price Elasticity of Demand)}: Một công ty SaaS tăng giá gói thuê bao từ \$40/tháng lên \$48/tháng ($+20\%$). Lượng người dùng trả phí giảm từ 5,000 xuống còn 4,200 người ($-16\%$). Tính hệ số co giãn của cầu theo giá ($E_d$). Quyết định tăng giá này có làm tăng hay giảm tổng doanh thu của công ty?
    \item \textbf{Bài 8 (Mô Hình Phân Bổ Tiếp Thị Đa Kênh -- Attribution Models)}: Một khách hàng thực hiện hành trình: Nhấp Quảng cáo Facebook (\$0) $\rightarrow$ Xem bài viết Blog Organic $\rightarrow$ Nhấp Google Search Ads $\rightarrow$ Mua hàng trị giá \$180. Hãy phân bổ doanh thu cho từng kênh theo 3 mô hình: First-Touch, Last-Touch, và Linear Attribution. Phân tích điểm yếu chí mạng của Last-Touch Attribution.
    \item \textbf{Bài 9 (Phân Khúc Khách Hàng RFM)}: Trình bày quy tắc xếp hạng Recency (R), Frequency (F), Monetary (M) từ 1 đến 5 điểm. Một khách hàng có điểm RFM là \texttt{1-5-5} (R=1, F=5, M=5) thuộc nhóm khách hàng nào? Doanh nghiệp cần hành động chiến lược gì ngay lập tức đối với nhóm này?
    \item \textbf{Bài 10 (Kinh Tế Học Đơn Vị -- Unit Economics Trong Giao Đồ Ăn)}: Một đơn hàng trà sữa có giá trị đơn (GTV) là 100,000 VND. Nguồn thu: Phí giao hàng thu từ khách: 15,000 VND; Hoa hồng chiết khấu từ quán ăn: 20\% GTV. Chi phí phát sinh: Trả tiền shipper: 18,000 VND; Voucher khuyến mãi công ty tài trợ: 8,000 VND; Phí xử lý cổng thanh toán: 1.5\% GTV; Chi phí hỗ trợ khách hàng và vận hành: 3,000 VND. Tính Contribution Margin (Biên lợi nhuận đóng góp) trên đơn hàng này. Doanh nghiệp đang lãi hay lỗ trên từng đơn?
    \item \textbf{Bài 11 (Vòng Quay Tồn Kho \& Điểm Đặt Hàng Lại -- DSI \& ROP)}: Giá vốn hàng bán (COGS) hàng năm là \$1,440,000. Giá trị hàng tồn kho trung bình là \$120,000. Doanh số tiêu thụ trung bình mỗi ngày là 40 sản phẩm. Thời gian nhà cung cấp giao hàng (Lead Time) là 10 ngày. Lượng hàng tồn kho an toàn (Safety Stock) duy trì ở mức 150 sản phẩm. Hãy tính Số ngày tồn kho bình quân (DSI) và Điểm đặt hàng lại (Reorder Point -- ROP).
    \item \textbf{Bài 12 (Net Promoter Score -- NPS)}: Khảo sát 2,000 khách hàng về mức độ sẵn sàng giới thiệu sản phẩm: 1,100 người chấm 9--10 điểm; 600 người chấm 7--8 điểm; 300 người chấm 0--6 điểm. Hãy tính chỉ số NPS. Vì sao việc nhiều nhà quản lý loại bỏ nhóm Passives (7--8 điểm) ra khỏi phân tích là một sai lầm nguy hiểm?
    \item \textbf{Bài 13 (Tỷ Lệ Triệt Tiêu Doanh Số -- Cannibalization Rate)}: Hãng ra mắt dòng tai nghe không dây mới \textit{"AirBuds Pro"} với giá \$200, bán được 10,000 chiếc trong tháng đầu. Tuy nhiên, doanh số của dòng tai nghe cũ \textit{"AirBuds Basic"} (giá \$120) bị giảm từ 15,000 chiếc xuống còn 9,000 chiếc. Tính tỷ lệ triệt tiêu (Cannibalization Rate) và cho biết quyết định ra mắt sản phẩm mới có làm tăng tổng doanh thu không?
    \item \textbf{Bài 14 (Tỷ Lệ Khiếu Nại Bồi Hoàn -- Chargeback Rate)}: Một cổng thương mại điện tử xử lý 50,000 giao dịch trong tháng 3. Đội phòng chống gian lận ghi nhận 450 giao dịch bị chủ thẻ khiếu nại lừa đảo và ngân hàng thực hiện hoàn tiền (Chargeback). Tính Chargeback Rate. So sánh với ngưỡng giới hạn 1.0\% (hoặc 0.9\%) của các tổ chức thẻ quốc tế (Visa/Mastercard) và nêu hậu quả nếu vượt ngưỡng.
    \item \textbf{Bài 15 (Phân Tích Giỏ Hàng -- Market Basket Analysis)}: Trong 5,000 giao dịch tại siêu thị: Có 1,000 hóa đơn chứa Bánh mì, 800 hóa đơn chứa Sữa tươi, và 400 hóa đơn chứa CẢ Bánh mì và Sữa tươi. Hãy tính Support, Confidence và Lift của luật kết hợp: $\{ \text{Bánh mì} \} \rightarrow \{ \text{Sữa tươi} \}$. Nhận xét xem hai sản phẩm này có mối quan hệ tương hỗ thực sự hay chỉ xuất hiện ngẫu nhiên.
    \item \textbf{Bài 16 (SaaS Net Revenue Retention -- NRR \& Quick Ratio)}: Đầu tháng, công ty có Doanh thu định kỳ tháng (Beginning MRR) là \$200,000. Trong tháng diễn ra các biến động:
    \begin{itemize}
        \item Khách hàng cũ nâng cấp gói (Expansion MRR): +\$24,000
        \item Khách hàng cũ hạ gói (Contraction MRR): -\$6,000
        \item Khách hàng hủy hợp đồng hoàn toàn (Churn MRR): -\$10,000
        \item Khách hàng hoàn toàn mới (New MRR): +\$30,000
    \end{itemize}
    Tính Tỷ lệ duy trì doanh thu thuần (NRR) và Chỉ số tăng trưởng nhanh (SaaS Quick Ratio). Doanh nghiệp có đạt tiêu chuẩn tăng trưởng bền vững của các quỹ đầu tư mạo hiểm (NRR $> 100\%$) không?
    \item \textbf{Bài 17 (Chỉ Số Chăm Sóc Khách Hàng -- CSAT vs FCR)}: Nhóm hỗ trợ khách hàng đạt Thời gian phản hồi lần đầu (First Response Time -- FRT) cực nhanh là 1.5 phút (vượt KPI đề ra là 5 phút), nhưng điểm hài lòng khách hàng (CSAT) lại giảm từ 4.6 xuống 3.7 sao. Khi phân tích sâu, chỉ số First Contact Resolution (FCR -- Giải quyết triệt để vấn đề ngay lần liên hệ đầu) chỉ đạt 42\%. Phân tích xung đột lợi ích giữa chỉ số đo lường tốc độ (Speed Metric) và chất lượng giải quyết (Quality Metric).
    \item \textbf{Bài 18 (Tác Động Của Tỷ Lệ Hoàn Hàng E-Commerce)}: Một cửa hàng thời trang trực tuyến bán 8,000 đơn hàng/tháng, giá bán bình quân \$60/đơn, giá vốn hàng bán \$25/đơn. Chi phí vận chuyển chiều đi do shop chịu là \$4/đơn. Tỷ lệ đổi trả hàng (Return Rate) là 20\%. Với mỗi đơn đổi trả, shop mất thêm \$5 chi phí vận chuyển hoàn về và kiểm định lại kho hàng. Tính lợi nhuận gộp thực tế của shop sau khi đã trừ toàn bộ chi phí đổi trả.
    \item \textbf{Bài 19 (Hoạch Định Năng Suất -- Capacity Utilization Trong Vận Tải)}: Một đội xe vận tải có 50 xe tải. Mỗi xe có khả năng chạy tối đa 26 ngày/tháng, bình quân chở được 10 tấn hàng/chuyến (mỗi ngày chạy 1 chuyến). Trong tháng, toàn đội xe đã thực hiện 1,040 chuyến xe và vận chuyển tổng cộng 8,840 tấn hàng. Tính Tỷ lệ sử dụng số chuyến (Trip Utilization Rate) và Tỷ lệ lấp đầy tải trọng bình quân (Capacity Load Factor).
    \item \textbf{Bài 20 (Thiết Kế Cây Chỉ Số Bắc Đẩu -- North Star Metric)}: Đối với nền tảng gọi xe công nghệ (Ride-hailing platform kết nối Hành khách và Tài xế), giải thích tại sao chọn \textit{"Số lượt tải ứng dụng (App Installs)"} hoặc \textit{"Tổng Doanh Thu (GMV)"} làm North Star Metric là một sai lầm ngắn hạn. Hãy đề xuất 1 North Star Metric chuẩn xác và 4 Input Metrics thuộc các chiều: Cung (Supply), Cầu (Demand), Chất lượng khớp lệnh (Matching Quality), và Trải nghiệm người dùng (UX).
\end{enumerate}
\end{baitapbox}

\begin{dapanbox}[Đáp Án \& Lời Giải Ngắn Gọn Cho 20 Bài Toán Nghiệp Vụ]
\begin{itemize}
    \item \textbf{Bài 1}: Doanh thu năm/khách = $85 \times 3.2 = \$272$. Lợi nhuận gộp năm = $272 \times 40\% = \$108.8$. $CLV = 108.8 / 0.20 = \textbf{\$544}$. Nếu Churn Rate giảm xuống $15\%$, $CLV = 108.8 / 0.15 = \textbf{\$725.33}$ (tăng trưởng tới \textbf{33.3\%} giá trị trọn đời).
    \item \textbf{Bài 2}: $\text{Blended CAC} = 120000 / 3000 = \textbf{\$40}$. $\text{Paid CAC} = 120000 / 2000 = \textbf{\$60}$. Lợi nhuận gộp hàng tháng/khách = $25 \times 80\% = \$20$. Thời gian hoàn vốn đối với khách Paid = $60 / 20 = \textbf{3.0 tháng}$. Đây là chỉ số rất tốt (chuẩn thị trường dưới 12 tháng).
    \item \textbf{Bài 3}: Retention Rates: M0 = 100\%, M1 = 65\%, M2 = 52\%, M3 = 48\%, M4 = 46\%. Đường cong giảm mạnh ở tháng đầu (mất 35\% người dùng do Onboarding hoặc chưa thấy giá trị cốt lõi), sau đó phẳng dần từ M2 đến M4 (giảm nhẹ 2--4\%/tháng), chứng minh sản phẩm đã đạt Product-Market Fit với tệp khách hàng trung thành ở mức $\sim 46\%$.
    \item \textbf{Bài 4}: $\text{End-to-End CR} = 3600 / 200000 = \textbf{1.8\%}$. Drop-off rates từng bước: Xem $\rightarrow$ Giỏ hàng: 75\%; Giỏ hàng $\rightarrow$ Bấm TT: 70\%; Bấm TT $\rightarrow$ Điền thông tin: 20\%; Điền thông tin $\rightarrow$ TT thành công: 70\%. Hai điểm rơi nghiêm trọng nhất là Xem $\rightarrow$ Thêm giỏ hàng (75\%) và Điền thông tin $\rightarrow$ Hoàn tất thanh toán (70\% do lỗi cổng thanh toán, phí ship bất ngờ, hoặc bắt buộc đăng ký phức tạp).
    \item \textbf{Bài 5}: Campaign A: Tổng CR = $(20 + 320)/(1000 + 4000) = 340/5000 = \textbf{6.8\%}$. Campaign B: Tổng CR = $(100 + 85)/(4000 + 1000) = 185/5000 = \textbf{3.7\%}$. Nghịch lý xảy ra vì Campaign A đổ tới 80\% traffic vào Google (kênh có CR cơ bản cao $8\%$), trong khi Campaign B đổ tới 80\% traffic vào Facebook (kênh có CR cơ bản thấp $2.5\%$). Biến ẩn (Confounding Variable) chính là tỷ trọng phân bổ nền tảng.
    \item \textbf{Bài 6}: Gọi tỷ lệ thay đổi của CR là $x$. Ta có phương trình chỉ số: $(1 - 0.12) = (1 + 0.10) \times (1 + x) \times (1 + 0.05) \Leftrightarrow 0.88 = 1.155 \times (1 + x) \Rightarrow 1 + x = 0.88 / 1.155 \approx 0.7619$. Như vậy, Conversion Rate đã \textbf{giảm sút tới $23.81\%$}. Cần rà soát ngay lập tức luồng checkout, khuyến mãi, hoặc chất lượng nguồn traffic mới kéo về.
    \item \textbf{Bài 7}: $\% \Delta P = +20\%$. $\% \Delta Q = (4200 - 5000)/5000 = -16\%$. $E_d = -16\% / 20\% = \textbf{-0.8}$. Vì $|E_d| < 1$, cầu co giãn ít (Inelastic Demand). Doanh thu cũ: $5000 \times 40 = \$200,000$; Doanh thu mới: $4200 \times 48 = \$201,600$ (tăng \$1,600). Quyết định tăng giá làm tăng tổng doanh thu.
    \item \textbf{Bài 8}: First-Touch: Facebook hưởng 100\% (\$180); Last-Touch: Google Ads hưởng 100\% (\$180); Linear: Mỗi kênh (FB, Blog, Google) hưởng 33.33\% (\$60). Điểm yếu của Last-Touch: Bỏ qua hoàn toàn vai trò xây dựng nhận thức thương hiệu (Awareness) của các kênh đầu phễu, dẫn đến quyết định sai lầm khi cắt ngân sách Facebook khiến nguồn khách mới cạn kiệt.
    \item \textbf{Bài 9}: Nhóm \texttt{1-5-5} là nhóm \textbf{"At Risk / Cannot Lose Them"} (Khách VIP có nguy cơ rời bỏ): Đã từng mua sắm cực kỳ thường xuyên và chi tiêu rất lớn, nhưng đã lâu không quay lại. Hành động chiến lược: Kích hoạt chiến dịch Re-engagement cá nhân hóa ngay, gọi điện chăm sóc trực tiếp từ quản lý tài khoản (Account Manager), hoặc tặng ưu đãi độc quyền cao cấp.
    \item \textbf{Bài 10}: Doanh thu gộp đơn hàng = $15000 + (100000 \times 20\%) = 35,000$ VND. Tổng chi phí trực tiếp trên đơn = $18000 + 8000 + (100000 \times 1.5\%) + 3000 = 18000 + 8000 + 1500 + 3000 = 30,500$ VND. $\text{Contribution Margin} = 35000 - 30500 = \textbf{+4,500 VND}$ (đạt tỷ suất lợi nhuận đóng góp $4.5\%$). Đơn hàng có lãi biên dương.
    \item \textbf{Bài 11}: Vòng quay tồn kho = $\text{COGS} / \text{Inventory} = 1440000 / 120000 = 12$ vòng/năm. $DSI = 365 / 12 \approx \textbf{30.42 ngày}$. Điểm đặt hàng lại: $ROP = (\text{Lead Time} \times \text{Daily Demand}) + \text{Safety Stock} = (10 \times 40) + 150 = \textbf{550 sản phẩm}$. Khi kho còn 550 sản phẩm, hệ thống phải tự động tạo đơn nhập hàng.
    \item \textbf{Bài 12}: $\% \text{Promoters} = 1100 / 2000 = 55\%$. $\% \text{Detractors} = 300 / 2000 = 15\%$. $NPS = 55\% - 15\% = \textbf{+40}$. Bỏ qua Passives là sai lầm vì nhóm này chiếm tới 30\% khách hàng; họ có mức độ trung thành thấp, rất dễ bị đối thủ lôi kéo chỉ bằng một chương trình giảm giá nhỏ nếu doanh nghiệp không kịp thời nâng cao trải nghiệm của họ.
    \item \textbf{Bài 13}: Lượng máy Basic sụt giảm: $15000 - 9000 = 6,000$ chiếc. $\text{Cannibalization Rate} = 6000 / 10000 = \textbf{60\%}$ (cứ 10 chiếc Pro bán ra thì có 6 chiếc lấy mất thị phần của Basic). Doanh thu cũ: $15000 \times 120 = \$1,800,000$. Doanh thu mới: $(9000 \times 120) + (10000 \times 200) = 1,080,000 + 2,000,000 = \$3,080,000$ (tăng \$1,280,000). Quyết định ra mắt sản phẩm mới cực kỳ thành công.
    \item \textbf{Bài 14}: $\text{Chargeback Rate} = 450 / 50000 = \textbf{0.90\%}$. Tỷ lệ này đã chạm sát ngưỡng cảnh báo đỏ của Visa/Mastercard (thường là 0.9\% -- 1.0\%). Doanh nghiệp đối mặt nguy cơ: Bị xếp vào chương trình giám sát rủi ro cao (Monitoring Program), phạt phí xử lý giao dịch tăng vọt từ \$15 lên \$50 mỗi đơn bồi hoàn, và nguy cơ nặng nhất là bị tước quyền thanh toán thẻ. Cần bật ngay xác thực 3D-Secure 2.0 và quy tắc kiểm duyệt địa chỉ IP/CVV.
    \item \textbf{Bài 15}: $\text{Support}(\text{Bánh mì} \rightarrow \text{Sữa}) = 400 / 5000 = \textbf{8.0\%}$. $\text{Confidence} = 400 / 1000 = \textbf{40.0\%}$ (40\% người mua Bánh mì sẽ mua Sữa). Tần suất nền của Sữa: $P(\text{Sữa}) = 800 / 5000 = 16\%$. $\text{Lift} = 40\% / 16\% = \textbf{2.5}$. Vì $\text{Lift} = 2.5 > 1$, sự xuất hiện của Bánh mì làm tăng khả năng mua Sữa tươi lên gấp 2.5 lần so với ngẫu nhiên $\rightarrow$ Rất nên tạo combo hoặc đặt gần quầy kệ.
    \item \textbf{Bài 16}: $NRR = \frac{200000 + 24000 - 6000 - 10000}{200000} \times 100\% = \frac{208000}{200000} \times 100\% = \textbf{104.0\%}$. Doanh nghiệp đạt chuẩn tăng trưởng bền vững (vượt mốc 100\%, tức là ngay cả khi không có khách hàng mới, doanh thu từ khách hàng cũ vẫn tự tăng trưởng 4\%). $\text{Quick Ratio} = \frac{\text{New MRR} + \text{Expansion MRR}}{\text{Churn MRR} + \text{Contraction MRR}} = \frac{30000 + 24000}{10000 + 6000} = \frac{54000}{16000} = \textbf{3.375}$ (rất khỏe mạnh, vì $> 2.0$).
    \item \textbf{Bài 17}: Khi nhân viên hỗ trợ bị áp lực KPI "Phản hồi dưới 2 phút", họ có xu hướng gửi câu trả lời đóng khung mẫu (canned responses) nhanh chóng để đóng ticket mà không thực sự lắng nghe hoặc giải quyết triệt để vấn đề gốc rễ. Khách hàng phải nhắn tin lại nhiều lần khiến CSAT tụt dốc. Cần cân bằng KPI bằng cách đưa FCR làm chỉ số bảo vệ (Guardrail Metric) và chỉ tính đạt KPI khi CSAT $\ge 4.5$.
    \item \textbf{Bài 18}: Số đơn hoàn trả: $8000 \times 20\% = 1,600$ đơn. Số đơn thành công thực nhận: $8000 - 1600 = 6,400$ đơn. Doanh thu thực nhận: $6400 \times 60 = \$384,000$. Giá vốn COGS hàng đã bán: $6400 \times 25 = \$160,000$. Phí ship chiều đi của toàn bộ 8,000 đơn ban đầu: $8000 \times 4 = \$32,000$. Chi phí hoàn hàng phát sinh: $1600 \times 5 = \$8,000$. Tổng lợi nhuận gộp sau hoàn hàng: $384000 - 160000 - 32000 - 8000 = \textbf{\$184,000}$. Tỷ suất lợi nhuận gộp thực tế: $184000 / 384000 \approx \textbf{47.9\%}$.
    \item \textbf{Bài 19}: Năng lực vận chuyển tối đa về số chuyến: $50 \times 26 = 1,300$ chuyến. $\text{Trip Utilization Rate} = 1040 / 1300 = \textbf{80.0\%}$. Khối lượng hàng tối đa có thể chở trên 1,040 chuyến đã chạy: $1040 \times 10 = 10,400$ tấn. $\text{Capacity Load Factor} = 8840 / 10400 = \textbf{85.0\%}$. Hiệu suất toàn diện của đội xe: $80\% \times 85\% = 68.0\%$.
    \item \textbf{Bài 20}: Lượt tải app và GMV là chỉ số bề nổi (Vanity Metrics): Tải app nhiều nhưng gỡ app ngay không tạo ra giá trị; GMV cao có thể do đốt voucher trợ giá thiếu bền vững. \textbf{North Star Metric}: \textit{"Số chuyến xe hoàn thành thành công hàng tuần (Weekly Completed Trips)"}. 4 Input Metrics cốt lõi: (1) Cung: Tỷ lệ tài xế online giờ cao điểm (Driver Active Hours); (2) Cầu: Tỷ lệ mở ứng dụng và tìm chuyến (Search-to-Request Conversion); (3) Chất lượng khớp lệnh: Thời gian chờ xe trung bình (ETA dưới 4 phút) và Tỷ lệ hủy chuyến của tài xế (Driver Cancellation Rate $< 3\%$); (4) Trải nghiệm: Điểm đánh giá chuyến đi trung bình $\ge 4.8$ sao.
\end{itemize}
\end{dapanbox}
